\documentclass{ximera}
%verbeterd 21 augustus 2026
\title{Aantrekking van ongeladen materie}
\author{Daan Lipkens}

\begin{document}
\begin{abstract}
    Het mechanisme achter het aantrekken van ongeladen deeltjes en hun toepassingen.
\end{abstract}
\maketitle


%------------------------------ AANTREKKING ONGELADEN MATERIE
\section{Aantrekking van ongeladen materie}

Een proces dat analoog is aan inductie in geleiders doet zich eveneens voor in isolatoren. Dit fenomeen wordt \textbf{polarisatie} genoemd. In een neutrale molecule valt het zwaartepunt van de positieve ladingen normaal gezien samen met dat van de negatieve ladingen. Wanneer een isolator zich echter in de nabijheid van een geladen voorwerp bevindt, kunnen deze zwaartepunten licht verschuiven. Hierdoor ontstaat binnen elke molecule een kleine ladingsscheiding, met als gevolg dat aan de ene zijde van de molecule een nettopositieve en aan de andere zijde een nettonegatieve lading optreedt. Deze microscopische herschikkingen leiden macroscopisch tot de vorming van een ladingslaag aan het oppervlak van de isolator, zoals weergegeven in figuur \ref{fig:scheiden van lading}a.

%---------------------------
% Figuur: ongeladen deeltjes
\begin{figure}[ht!]
    \begin{center}
        \begin{tabular}{cc}
            \includegraphics[width=4.5cm]{pic5/elektriciteit/coulomb/ballon en muur.png} &
            \includegraphics[width=4.5cm]{pic5/elektriciteit/coulomb/aantrekken papier.jpg} \\
            (a) & (b)
        \end{tabular}
    \end{center}
    \caption{(a) Het geladen voorwerp aan de linkerkant wekt een ladingsverdeling op aan het oppervlak van een isolator door de heroriëntatie van de ladingen in de moleculen. (b) Een geladen kam trekt stukjes papier aan omdat de ladingen in de moleculen van het papier zich herschikken.}
    \label{fig:scheiden van lading}
\end{figure}
%------------------------------


Dit principe verklaart waarom een geladen kam, nadat deze door het haar is gehaald, kleine stukjes papier aantrekt (figuur \ref{fig:scheiden van lading}b). Door de wrijving met het haar wordt de kam geladen. De nabijheid van de geladen kam veroorzaakt vervolgens een verschuiving van de ladingen in de moleculen van het papier. De tegengestelde ladingen bevinden zich nu dichter bij de kam, terwijl de gelijksoortige ladingen juist verder weggeduwd worden. Omdat de elektrostatische kracht volgens vergelijking \ref{eq:coulombkracht} omgekeerd evenredig is met het kwadraat van de afstand ($F \sim \frac{1}{r^{2}}$), is de aantrekkingskracht sterker dan de afstotingskracht. Het resultaat is een netto aantrekkende kracht, waardoor de stukjes papier naar de kam toe bewegen.

\section{Toepassingen}
Deze toepassingen zijn geen leerstof, maar kunnen wel gebruikt worden om het begrip van de leerstof te vergroten en hun nut te zien in het dagelijks leven.
% TOEPASSING: SWIFFER
\begin{example}[foldable=true,title={Toepassing: Swiffer}]\nl
    Een \textbf{Swifferdoekje} werkt op basis van elektrostatische aantrekking. Tijdens het productieproces wordt het doekje zodanig behandeld dat het gemakkelijk een elektrostatische lading opneemt. Wanneer je ermee over de vloer of meubels veegt, wekt de wrijving een lading op in de vezels van het doekje.

    Stofdeeltjes en haren in de omgeving zijn elektrisch neutraal, maar de lading in het doekje zorgt voor een kleine polarisatie in die deeltjes. Hierdoor worden ze aangetrokken door het doekje en blijven ze eraan kleven. Zo verwijdert een Swiffer veel meer stof dan een gewoon droog doekje, zonder gebruik van water of schoonmaakmiddelen.
    
    \begin{center}
        \includegraphics[width=5cm]{pic5/elektriciteit/coulomb/swiffer.jpg}
        \captionof{figure}{Swiffer stofveegdoekje.}
        \label{fig:swiffer}
    \end{center}
\end{example}

%---------------
% TOEPASSING SPRAY
\begin{example}[foldable=true,title={Toepassing: Poedercoating en elektrostatisch spuiten}]\nl
    \begin{itemize}
        \item \textbf{Poedercoating} is een milieuvriendelijk alternatief voor het verven van bijvoorbeeld ijzer of aluminium. Hierbij wordt het poeder elektrostatisch geladen, terwijl het te bekleden materiaal een tegengestelde lading krijgt. Door de elektrostatische aantrekking hecht het poeder zich gelijkmatig aan het oppervlak en bereikt het zelfs de kleinste hoekjes. Vervolgens wordt het voorwerp in een oven verhit, waardoor het poeder smelt en een beschermende laag vormt tegen corrosie. Op die manier zijn er geen schadelijke oplosmiddelen nodig.
        \item \textbf{Elektrostatisch spuiten van gewassen} vermindert het gebruik van spuitvloeistof met ongeveer de helft. Door de vloeistofdruppels elektrostatisch te laden, hechten ze zich gelijkmatig aan de gewassen, waardoor één enkele bespuiting een egaal resultaat oplevert. Dit effect is te zien op figuur \ref{fig:spray}.
    \end{itemize}
    %------------------------------------
    % Figuur spray
    \begin{center}
        \includegraphics[width=4.5cm]{pic5/elektriciteit/coulomb/besproeien.jpg}
        \captionof{figure}{Bij het besproeien van gewassen kan men de sproeistof eerst elektrostatisch laden, waardoor minder sproeistof nodig is.}
        \label{fig:spray}
    \end{center}
    %------------------------------------------------------
    % Video pesticide spuiten
    \begin{onlineOnly}
        \youtube{jPliFIgEWYU}
    \end{onlineOnly}
    \pdfOnly{
        \begin{center}
            \mylink{https://youtu.be/jPliFIgEWYU}{Link naar video over pesticide spuiten op gewassen met elektrostatica.}
        \end{center}
    }
\end{example}
%------------------------------


%---------------------------
% TOEPASSING: STOFFILTRATIE
\begin{example}[foldable=true,title={Toepassing: Stoffiltratie}]\nl
    In fabrieken, elektriciteitscentrales en andere industriële installaties ontstaan bij verbrandingsprocessen vaak rookgassen die stofdeeltjes bevatten. Om de uitstoot van deze deeltjes zoveel mogelijk te beperken, worden de stofdeeltjes elektrostatisch geladen en vervolgens langs een rooster met tegengestelde lading geleid. Daar worden ze door elektrostatische aantrekking vastgehouden en zo uit de lucht verwijderd.
    
    Dit principe van \textbf{elektrostatische stoffiltratie} wordt niet alleen toegepast in grote installaties, maar ook in luchtfilters voor woningen en kantoren. Een dergelijk filter is te zien op figuur \ref{fig:stoffiltratie}.
    %------------------------------------
    % Figuren stoffiltratie
    \begin{center}
        \includegraphics[width=4.5cm]{pic5/elektriciteit/coulomb/stofvanger.jpg}
        \captionof{figure}{In fabrieken worden neutrale rookdeeltjes elektrostatisch aangetrokken door geladen roosters en zo uit de lucht gefilterd.}
        \label{fig:stoffiltratie}
    \end{center}
    %------------------------------------------------------
\end{example}

%---------------
% TOEPASSING kopieerapparaat
\begin{example}[foldable=true,title={Toepassing: Laserprinter en kopieerapparaat}]\nl
    Een \textbf{laserprinter} of \textbf{kopieerapparaat} maakt gebruik van elektrostatische ladingen om tekst en afbeeldingen op papier over te brengen.

    Eerst wordt een cilindervormige drum elektrisch geladen met behulp van een hoogspanningsbron. Vervolgens tekent een laserstraal of ledlicht de gewenste tekst of afbeelding op de drum. Daar waar de laser invalt, wordt de lading lokaal geneutraliseerd. Zo ontstaat er een patroon van geladen en ongeladen gebieden dat overeenkomt met de af te drukken afbeelding.

    Daarna wordt \textbf{tonerpoeder}, dat zelf elektrostatisch geladen is, op de drum gestrooid. Het poeder blijft enkel kleven op de plaatsen die nog geladen zijn. Het patroon van de toner vormt zo een onzichtbare ``kopie'' van de tekst of afbeelding op de drum.

    In de volgende stap rolt er een vel papier langs de drum. Het papier krijgt een tegengestelde lading en trekt daardoor het tonerpoeder van de drum aan. Ten slotte wordt het papier door verwarmde rollen geleid, waardoor de toner in het papier smelt en zich stevig vasthecht.

    Op deze manier kan een printer of kopieerapparaat snel en efficiënt scherpe afdrukken maken dankzij elektrostatische principes. (Zie ook figuur \ref{fig:printer}).
    
    %------------------------------------
    % Figuren printer
    \begin{center}
        \includegraphics[width=4.5cm]{pic5/elektriciteit/coulomb/printer.jpg}
        \captionof{figure}{In een printer wordt de toner met behulp van elektrostatische ladingen op de juiste plaatsen aangebracht.}
        \label{fig:printer}
    \end{center}
    %------------------------------------------------------
    % Video laserprinter
    \begin{onlineOnly}
        \youtube{Sy5PPaxaBCI}
    \end{onlineOnly}
    \pdfOnly{
        \begin{center}
            \mylink{https://youtu.be/Sy5PPaxaBCI}{Link naar video met uitleg over de werking van het kopieerapparaat.}
        \end{center}
    }
\end{example}

\end{document}